\section{Background and Related works}

\subsection{Diffusion Language Models}
Diffusion models have emerged as a powerful paradigm in generative modeling,
framing data generation as the inversion of a forward-noise process. They have achieved
remarkable success in continuous domains such as images and videos. More
recently, diffusion models have also advanced the natural language processing area.
DLMs extend diffusion to discrete sequences by redefining noise injection and
denoising.

DLMs generate text via an iterative unmasking process
over $T$ discrete denoising steps, gradually transforming a masked sequence into
the final output. Let $\mathcal V$ denote the vocabulary, and let
$x_t^l \in \mathcal V^l$ denote the sequence state of length $l$ at step $t$,
where $t = 0, \ldots, T$.

If the prompt is formulated as $p_0 = (x_1, \ldots, x_m)$, the goal is to generate
a full sequence $x = (x_1, \ldots, x_L)$ with the prompt tokens fixed. We define
the response positions as
\[
\mathcal R = \{m+1, \ldots, L\}.
\]

At inference time, DLMs start from the maximum-noise state by masking all response
tokens:
\[
x_i^{0} =
\begin{cases}
x_i, & i \le m, \\
[\mathrm{MASK}], & i \in \mathcal R,
\end{cases}
\]
where $[\mathrm{MASK}]$ denotes a special mask token. Prompt tokens are frozen
throughout the inference process.

Given a decreasing mask-ratio schedule $1=t_0>t_1>\cdots>t_K=0$, inference proceeds
for $K$ steps. At step $k$, a Transformer denoiser $f_\theta$ produces token
distributions
\[
\pi_i^{k} = \mathrm{Softmax}(f_\theta(x^k)_i), \quad i \in \mathcal R.
\]
Let $\mathcal M_k = \{ i \in \mathcal R \mid x_i^{k} = [\mathrm{MASK}] \}$. A subset
$\mathcal U_k \subseteq \mathcal M_k$ with
\[
|\mathcal U_k| = |\mathcal M_k| - \lceil t_{k+1}|\mathcal R| \rceil
\]
is unmasked by sampling from $\pi_i^{k}$, while all other positions remain
unchanged. After $K$ steps, all masks are removed and $x^K$ is returned.


\subsection{Efficient inference for Diffusion Models}
To balance the trade-off between inference speed and model performance, various techniques
have been proposed. One of the most accepted methods is block diffusion, which sampling
applies the origin reverse process within each block and the autoregressive sampling
across blocks. This approach activates the KV cachce strategies in DLLMs, which can 
significantly reduce the computational overhead during inference. Other techniques
includes \zz{literature summary}

\subsection{Sparse Attention}

In long-context language models, sparse attention mechanisms have been widely adopted to
reduce the quadratic complexity of standard attention. Various sparse attention patterns